% SGA 3, 第 VI_B 篇, 第 1 轮英文重构.
% 范围: 引理 5.6.3 全文, 含底本编者注 44, 以及随后的定理 5.3 的推导.
% 底本: Polo--Gille Exp6B-13oct24.pdf, 分篇第 33--34 页
% (合订本第 374--375 页; 原版页边印刷页码 355--356, 重版本页眉 363--364).
% 反例 5.6.4 始于分篇第 34 页 / 合订本第 375 页.
% OCR 与 Reinhold 英文 Markdown 仅为比对材料.
% 本单元无图.

\medskip
\phantomsection\label{lem:sga3-VIB-5.6.3}%
\noindent\textbf{引理 5.6.3.---}
\emph{设 \(X\) 是 \(S\)-概形. 下列条件等价:}
\begin{enumerate}
\item[\textup{\SGAInlineRefTarget{lem:sga3-VIB-5.6.3:item:i}(i)}]\phantomsection\label{cond:sga3-VIB-5.6.3-i}
  \emph{\(X\) 在 \(S\) 上分离.}

\item[\textup{\SGAInlineRefTarget{lem:sga3-VIB-5.6.3:item:ii}(ii)}]\phantomsection\label{cond:sga3-VIB-5.6.3-ii}
  \emph{对每个 \(S\)-概形 \(T\), 每个截面 \(\sigma:T\to X_T\) 都是闭浸入.}

\item[\textup{\SGAInlineRefTarget{lem:sga3-VIB-5.6.3:item:iii}(iii)}]\phantomsection\label{cond:sga3-VIB-5.6.3-iii}
  \emph{对每个既约 \(S\)-概形 \(T\),
  在 \(T\) 的某个稠密开集 \(U\) 上相同的两个 \(S\)-态射
  \(f_1,f_2:T\to X\) 相等.}

\item[\textup{\SGAInlineRefTarget{lem:sga3-VIB-5.6.3:item:iv}(iv)}]\phantomsection\label{cond:sga3-VIB-5.6.3-iv}
  \emph{对每个 \(s\in S\) 及 \(X_s\) 中的每对点 \(x_1,x_2\),
  存在态射}
  \[
    \phi:S''\longrightarrow S'\longrightarrow S
  \]
  \emph{及 \(X_{S''}\) 的一个在 \(S''\) 上分离的开子概形 \(V\), 使得:}
  \begin{enumerate}
  \item[\textup{(a)}]\phantomsection\label{cond:sga3-VIB-5.6.3-iv-a}
    \emph{\(S'\to S\) 开, \(S''\to S'\) 闭且满射,
    并且 \(\phi^{-1}(s)\ne\varnothing\);}
  \item[\textup{(b)}]\phantomsection\label{cond:sga3-VIB-5.6.3-iv-b}
    \emph{\(\{x_1,x_2\}\) 在 \(X_{S''}\) 中的逆像包含在 \(V\) 中.}
  \end{enumerate}

\item[\textup{(iv\(^{\prime}\))}]
  \phantomsection\label{cond:sga3-VIB-5.6.3-iv-prime}
  \emph{对每个 \(s\in S\), 每对点 \(x_1,x_2\in X_s\)
  都包含在 \(X\) 的一个在 \(S\) 上分离的开集 \(V\) 中.}%
  \footnote{\label{note:sga3-VIB-44}编者注: 我们简化了条件
  \textup{\SGARefLink{lem:sga3-VIB-5.6.3:item:iv}{(iv)}} 的表述,
  并增加条件 \textup{(iv\(^{\prime}\))}.
  还补出了 \textup{\SGARefLink{lem:sga3-VIB-5.6.3:item:i}{(i)}}\(\Rightarrow\)\textup{\SGARefLink{lem:sga3-VIB-5.6.3:item:iii}{(iii)}}
  的证明, 它用于
  \textup{\SGARefLink{lem:sga3-VIB-5.6.3:item:iv}{(iv)}}\(\Rightarrow\)\textup{\SGARefLink{lem:sga3-VIB-5.6.3:item:iii}{(iii)}}
  的证明. 此外注意, 若
  \[
    T=\operatorname{Spec}B,\qquad
    B=k[\varepsilon,x]/(\varepsilon^2,\varepsilon x)
  \]
  \textup{(\(k\) 是域)}, 且
  \[
    X=T_{\mathrm{red}}=\operatorname{Spec}k[x],
  \]
  则态射
  \[
    \phi_\lambda:T\longrightarrow X,\qquad
    x\longmapsto x+\lambda\varepsilon
    \quad(\lambda\in k),
  \]
  在稠密开集 \(\operatorname{Spec}B_x\) 上相同, 但并不相等.}
\end{enumerate}

\emph{证明.}
蕴含关系 \textup{\SGARefLink{lem:sga3-VIB-5.6.3:item:ii}{(ii)}}\(\Rightarrow\)\textup{\SGARefLink{lem:sga3-VIB-5.6.3:item:i}{(i)}}
显然 (取 \(T=X\), 并令 \(\sigma\) 为对角截面), 同样显然的还有
\[
  \textup{(i)}\Rightarrow\textup{(iv\(^{\prime}\))}
  \Rightarrow\textup{(iv)}.
\]
另一方面, 由 EGA~I, 5.4.6,
\(\textup{(i)}\Rightarrow\textup{(ii)}\).

\smallskip
\noindent\textup{\SGARefLink{lem:sga3-VIB-5.6.3:item:iii}{(iii)}}\(\Rightarrow\)\textup{\SGARefLink{lem:sga3-VIB-5.6.3:item:ii}{(ii)}}.
设 \(\sigma:T\to X_T\) 是 \(p:X_T\to T\) 的截面.
由 EGA~I, 5.3.13, \(\sigma\) 是浸入,
即从 \(T\) 到 \(X_T\) 的某个局部闭子概形 \(E\) 的同构.
为证明 \(E\) 闭, 可假设 \(T\) 和 \(E\) 既约.
设 \(\overline E\) 是 \(X_T\) 中以 \(E\) 的闭包为底空间的既约闭子概形,
从而 \(E\) 是 \(\overline E\) 的稠密开子概形. 浸入
\[
  i:\overline E\hookrightarrow X_T
\]
与 \(\sigma\circ p\circ i\) 在 \(E\) 上相同,
因而由假设 \textup{\SGARefLink{lem:sga3-VIB-5.6.3:item:iii}{(iii)}},
在 \(\overline E\) 上相同. 所以 \(\overline E\) 的每个点都属于
\(\sigma(T)=E\), 从而 \(\overline E=E\).

\smallskip
\noindent\textup{\SGARefLink{lem:sga3-VIB-5.6.3:item:i}{(i)}}\(\Rightarrow\)\textup{\SGARefLink{lem:sga3-VIB-5.6.3:item:iii}{(iii)}}.
假设 \(X\) 在 \(S\) 上分离, 并设 \(T\) 是既约 \(S\)-概形,
\(f_1,f_2:T\to X\) 是两个在 \(T\) 的稠密开集 \(U\) 上相同的 \(S\)-态射.
设
\[
  g:T\longrightarrow X\times_S X
\]
是以 \(f_1\) 和 \(f_2\) 为分量的态射. 由于
\[
  D=\Delta_{X/S}(X)
\]
闭, 它在 \(g\) 下的逆像是 \(T\) 的一个包含 \(U\) 的闭集,
因而等于 \(T\). 由于 \(T\) 既约, \(g\) 经过 \(D\) 分解
(参见 EGA~I, 5.2.2); 因此
\[
  f_1=p_1\circ g=p_2\circ g=f_2.
\]

\smallskip
\noindent\textup{\SGARefLink{lem:sga3-VIB-5.6.3:item:iv}{(iv)}}\(\Rightarrow\)\textup{\SGARefLink{lem:sga3-VIB-5.6.3:item:iii}{(iii)}}.
设 \(T\) 是既约 \(S\)-概形, \(f_1,f_2:T\to X\)
是两个在稠密开集 \(U\) 上相同的 \(S\)-态射.
由于 \(T\) 既约, 为证明 \(f_1=f_2\),
只须证明 \(f_1=f_2\) 在集合论意义下成立.
事实上, 假设已证明此点, 取 \(t\in T\),
设 \(V\) 是 \(X\) 中包含 \(f_1(t)=f_2(t)\) 的仿射开集,
\(W\) 是 \(t\) 的开邻域,
它是 \(V\) 在 \(f_1\) 和 \(f_2\) 共同的底连续映射下的逆像.
则态射
\[
  f_i|_W:W\longrightarrow V
\]
在 \(W\) 的稠密开集 \(U\cap W\) 上相同.
由于 \(V\) 分离且 \(W\) 既约, 可知 \(f_1|_W=f_2|_W\),
从而 \(f_1=f_2\).

现取 \(t\in T\), 设 \(s\) 是它在 \(S\) 中的像,
证明 \(X_s\) 的点
\[
  x_1=f_1(t),\qquad x_2=f_2(t)
\]
相等. 取
\[
  \phi:S''\longrightarrow S'\longrightarrow S
\]
及 \(X\times_S S''\) 的开集 \(V\),
如 \textup{\SGARefLink{lem:sga3-VIB-5.6.3:item:iv}{(iv)}} 中所述. 置
\[
  T'=T\times_S S',\qquad T''=T\times_S S'',
\]
并用
\[
  g:T''\longrightarrow T,\qquad
  f_i'':T''\longrightarrow X_{S''}\quad(i=1,2)
\]
表示基变换所得的态射.

由于 \(U\) 在 \(T\) 中稠密, 且 \(T'\to T\) 开,
\(U\) 在 \(T'\) 中的逆像 \(U'\) 在 \(T'\) 中稠密.
设 \(U''\) 是 \(U'\) 在 \(T''\) 中的逆像,
\(F\) 是 \(T''\) 的以 \(\overline{U''}\) 为底空间的既约子概形.
由于 \(T''\to T'\) 满射且闭, \(F\) 的像包含 \(U'\) 且闭,
因而等于 \(T'\). 所以
\[
  F\cap g^{-1}(t)
\]
包含一点 \(u\).

对 \(i=1,2\), 设 \(h_i\) 是 \(f_i''\) 在 \(F\) 上的限制.
则 \(h_i(u)\) 是 \(X_{S''}\) 中位于 \(f_i(t)=x_i\) 上的点,
因而属于 \(V\), 因为 \(V\) 包含 \(\{x_1,x_2\}\) 在 \(X_{S''}\) 中的逆像.

于是
\[
  W=h_1^{-1}(V)\cap h_2^{-1}(V)
\]
是 \(F\) 的包含 \(u\) 的开集, 且 \(S''\)-态射
\[
  h_i|_W:W\longrightarrow V
\]
在 \(W\) 的稠密开集 \(U''\cap W\) 上相同.
由于 \(V\) 在 \(S''\) 上分离, 有
\[
  h_1|_W=h_2|_W.
\]
因此 \(h_1(u)=h_2(u)\), 从而 \(f_1(t)=f_2(t)\).
这证明了 \(\textup{(iv)}\Rightarrow\textup{(iii)}\).

\medskip
\phantomsection\label{ded:sga3-VIB-theorem-5.3-from-5.6.2-5.6.3}%
\hyperref[thm:sga3-VIB-5.3]{定理 5.3} 现在由
\hyperref[cor:sga3-VIB-5.6.2]{5.6.2} 及
\hyperref[lem:sga3-VIB-5.6.3]{5.6.3} 的蕴含关系
\(\textup{(iv)}\Rightarrow\textup{(i)}\) 得到.
